\section*{Hoofdstuk \ref{ch:g9:earth-resources} --- Grondstoffen uit de Aarde}

\begin{solution}{exo:g9:earth-resources:1}
(a) \omterm{def:g9:earth-resources:aggregate}{Toeslagmateriaal}. (b) \omterm{def:g9:earth-resources:reserve}{Reserve}. (c) \omterm{def:g9:earth-resources:ore-deposit}{Ertsafzetting}. (d) \omterm{def:g9:earth-resources:ore-deposit}{Concentratiefactor}.
\end{solution}

\begin{solution}{exo:g9:earth-resources:2}
(a) In riviergrind. (b) In spleten, uit heet water. (c) In het \omterm{def:g7:volcanoes:magma}{magma}.
(d) Dit is de eerste stap van een afzetting in riviergrind: het vrijgemaakte
goud wordt daarna door de \omterm{def:g1:landscapes-around-us:water}{rivier} meegevoerd.
\end{solution}

\begin{solution}{exo:g9:earth-resources:3}
(b), (d), (a), (e), (c).
\end{solution}

\begin{solution}{exo:g9:earth-resources:4}
(a) Niet hernieuwbaar: het heeft miljoenen jaren nodig om te ontstaan.
(b) Hernieuwbaar: de \omterm{def:g4:clouds-rain-snow:precipitation}{regen} vult aan wat er wordt weggenomen. (c) Niet
hernieuwbaar: een \omterm{def:g9:earth-resources:ore-deposit}{ertsafzetting} ontstaat over heel lange tijden. (d) Niet
hernieuwbaar in dat tempo: de \omterm{def:g6:water-underground:water-table}{grondwaterspiegel} daalt en de putten vallen
uiteindelijk droog.
\end{solution}

\begin{solution}{exo:g9:earth-resources:5}
(a) $4 \div 0.007 \approx 571$, ongeveer $10^3$ (tussen $10^2$ en $10^3$,
dichter bij $10^3$). (b) $0.5 \div 0.0002 = 2500$, ongeveer $10^3$. (c) $20
\div 5.6 \approx 3.6$, ongeveer 4: tussen $10^0$ en $10^1$.
(d) IJzer komt in de korst al heel veel voor: een erts dat maar een paar
keer rijker is dan gewoon \omterm{def:g3:stones-pebbles-sand:rock}{gesteente} is genoeg.
\end{solution}

\begin{solution}{exo:g9:earth-resources:6}
(a) $19.3 \div 2.5 \approx 7.7$ keer. (b) Goud: $19.3 \times 0.01 =
\qty{0.193}{g}$; \omterm{def:g3:stones-pebbles-sand:rock}{gesteente}: $2.5 \times 0.01 = \qty{0.025}{g}$. (c) Als de
hele bedding beweegt, zakken de zware goudkorrels tussen het lichtere \omterm{def:g3:stones-pebbles-sand:sand}{zand}
en de \omterm{def:g3:stones-pebbles-sand:pebble}{kiezels} naar beneden, tot op het \omterm{def:g3:stones-pebbles-sand:rock}{vaste gesteente}, waar ze blijven
liggen.
\end{solution}

\begin{solution}{exo:g9:earth-resources:7}
(a) Goud, tin en chroom. (b) Tussen $10^1$ en $10^2$, dicht bij $10^2$.
(c) $3000 \div 3.7 \approx 800$ keer. (d) Chroom (of tin, of goud): hoe
hoger de \omterm{def:g9:earth-resources:ore-deposit}{concentratiefactor}, hoe zeldzamer de plaatsen waar de natuur het
metaal genoeg heeft geconcentreerd.
\end{solution}

\begin{solution}{exo:g9:earth-resources:8}
(a) $18 \times 365 = \qty{6570}{kg}$, ongeveer \qty{6.6}{t}. (b) $5 \times
10^{10}$ ton. (c) $25000 \div 18 \approx 1389$ dagen, bijna vier jaar.
(d) Het \omterm{def:g3:stones-pebbles-sand:sand}{zand} dat uit de \omterm{def:g1:landscapes-around-us:water}{rivier} wordt gehaald, gaat niet meer stroomafwaarts
naar de \omterm{def:g1:landscapes-around-us:water}{kust}; de golven blijven \omterm{def:g3:stones-pebbles-sand:sand}{zand} van het strand wegvoeren, maar er komt
minder bij om het te vervangen.
\end{solution}

\begin{solution}{exo:g9:earth-resources:9}
(a) Van piepkleine planten en dieren uit de \omterm{def:g2:water-on-earth:ocean}{zee}, zoals diatomeeën, die in
de modder van de zeebodem zijn begraven. (b) Tussen ongeveer 60 en
\qty{120}{\degreeCelsius}. (c) Ze zijn lichter dan water, dat de \omterm{def:g6:water-underground:porosity}{poriën}
van de \omterm{def:g3:stones-pebbles-sand:rock}{gesteenten} vult. (d) Een ondoorlatende laag, zoals schalie; ze
verzamelen zich eronder, vaak aan de top van een opwelving, gas boven
olie.
\end{solution}

\begin{solution}{exo:g9:earth-resources:10}
(a) $2.9 \times 10^7$ ton. (b) $2.9 \times 10^7 \div 100 = 2.9 \times
10^5$ ton per jaar, ruwweg. (c) Het erts is 100 keer zwaarder dan het
koper dat het bevat: $2.9 \times 10^7 \times 100 = 2.9 \times 10^9$ ton.
(d) Elke ton erts levert maar \qty{10}{kg} koper, dus moeten er miljarden
tonnen \omterm{def:g3:stones-pebbles-sand:rock}{gesteente} worden afgegraven.
\end{solution}

\begin{solution}{exo:g9:earth-resources:11}
(a) Bij 23~miljoen ton per jaar zou ze ongeveer $280 \div 23
\approx 12$ jaar hebben meegegaan, tot het begin van de jaren 1980; toch
werd er 697~miljoen ton gewonnen, veel meer dan die \omterm{def:g9:earth-resources:reserve}{reserve}.
(b) $980 \div 280 = 3.5$ keer. (c) $980 \div
23 \approx 43$ jaar. (d) Bijvoorbeeld: er worden nieuwe afzettingen
ontdekt; betere technieken of hogere prijzen maken armere afzettingen het
winnen waard.
\end{solution}

\begin{solution}{exo:g9:earth-resources:12}
(a) $5 \times 10^1$ jaar. (b) $2 \times 10^3 \div 50 = 40$ keer. (c) $3
\times 10^8 \div (5 \times 10^1) = 6 \times 10^6$ keer. (d) Het ontstaan
van de steenkool duurde miljoenen keren langer dan wij nodig hebben om
haar te gebruiken: de natuur vult haar op onze tijdschaal niet aan.
\end{solution}

\begin{solution}{pb:g9:earth-resources:1}
\textbf{1.} $100 \times 1000 = \num{100000}$~kg $= 1 \times 10^5$~kg.

\textbf{2.} $2000 \times 300 = \num{600000}$~t $= 6 \times 10^5$~t.

\textbf{3.} $6 \times 10^5 \div 10 = 6 \times 10^4$~t per jaar.

\textbf{4.} $0.40 \times 6 \times 10^4 = 2.4 \times 10^4$~t \omterm{def:g3:stones-pebbles-sand:sand}{zand} per
jaar.

\textbf{5.} $1200 \times 500 \times 5 = 3 \times 10^6$~m$^3$.

\textbf{6.} Twee derde ervan: $2 \times 10^6$~m$^3$.

\textbf{7.} $2 \times 10^6 \times 1.6 = 3.2 \times 10^6$~t.

\textbf{8.} $\num{24000} + \num{136000} = \num{160000}$~t $= 1.6 \times
10^5$~t per jaar.

\textbf{9.} $\num{160000} \div 2000 = 80$ keer meer.

\textbf{10.} Nee: de groeve wordt 80 keer sneller leeggehaald dan de \omterm{def:g1:landscapes-around-us:water}{rivier}
het \omterm{def:g3:stones-pebbles-sand:sand}{zand} aanvult. Op de schaal van de stad is dit \omterm{def:g3:stones-pebbles-sand:sand}{zand} een eindige
grondstof.

\textbf{11.} $\num{24000} \div 4 = \num{6000}$~t \omterm{def:g3:stones-pebbles-sand:sand}{zand} per jaar.

\textbf{12.} $3.2 \times 10^6 \div (1.6 \times 10^5) = 20$: \textbf{20 jaar}.
\end{solution}
